\section{交叉熵训练其实在最小化 KL 差异}
\label{sec:13-Machine-Learning/11-Information-Theoretic-Learning/01}

你训练一个三分类模型。对某张输入图片，模型输出 softmax 概率 \(q = [0.10, 0.85, 0.05]\)，真实标签是第二类。你拿到的单样本损失是 \(-\log(0.85) \approx 0.163\)。如果模型错误地自信——\(q = [0.90, 0.05, 0.05]\)——损失变成 \(-\log(0.05) \approx 2.996\)。同样的分类错误，自信的程度差了 18 倍的惩罚。

这个对数惩罚从哪里来？为什么分类器训练几乎千篇一律地使用 cross-entropy，而不是更直观的 squared error——比如直接惩罚 \((1 - 0.85)^2 = 0.0225\)？答案不在损失函数的代数形式里，而在信息论对``用错误模型描述数据''给出了一个精确的代价度量。

\subsection{编码代价就是分布不匹配的度量}

设真实标签的分布是 \(p(y)\)——在监督学习中，给定输入 \(x\) 后标签的条件分布。模型输出的预测分布是 \(q(y)\)。

信息论里有一个操作性问题：如果你按照 \(q\) 设计了一套最优编码方案（给概率高的标签分配短码，概率低的分配长码），却用它来编码实际服从 \(p\) 的标签，平均每个标签需要多少比特？

答案是 cross-entropy：

\[
H(p, q) = -\sum_{y} p(y) \log q(y).
\]

它衡量的是``用模型 \(q\) 的码本描述真实来源 \(p\)''所付出的平均编码代价。这个代价可以拆成两部分：

\[
H(p, q)
= -\sum_y p(y)\log p(y)
- \sum_y p(y)\log\frac{q(y)}{p(y)}
= H(p) + \KL(p\|q).
\]

第一项 \(H(p)\) 是来源自身的不确定性——即使你知道了真实分布 \(p\) 并设计了最优编码，每个标签仍至少需要这么多比特。这是信息论的下限，无法通过改进模型来降低。

第二项 \(\KL(p\|q)\) 是 KL 散度——因为你使用了 \(q\) 而不是真实 \(p\) 来设计编码，额外多付的比特数。KL 非负，且只有当 \(q = p\) 时才归零。

对于固定的数据生成分布 \(p\)，\(H(p)\) 与模型参数无关。因此最小化期望 cross-entropy 精确等价于最小化 \(\KL(p\|q)\)。最优可达到的损失值是 \(H(p)\)——真实分布的熵，而不是零。一个完美模型在 cross-entropy 下的损失不是零，而是数据本身固有的不确定性。

\subsection{单标签训练如何与这个框架对接}

监督学习的训练数据每次只给出一个具体的标签 \(y_i\)，不是完整的分布 \(p(y\mid x_i)\)。one-hot 标签 \(\delta_{y_i}\) 与模型预测 \(q_\theta(\cdot\mid x_i)\) 之间的 cross-entropy 退化为：

\[
-\sum_{k=1}^K \ind\{y_i = k\} \log q_\theta(k\mid x_i)
= -\log q_\theta(y_i\mid x_i).
\]

这就是你在训练日志里看到的那个数。对全体训练样本取平均：

\[
-\frac{1}{n}\sum_{i=1}^n \log q_\theta(y_i\mid x_i).
\]

它同时是三个东西：categorical negative log-likelihood——给定参数后观测到这些标签的对数似然的负数；期望 cross-entropy \(H(p, q)\) 的样本估计——用经验分布替代了真实的 \(p\)；以及标准分类训练的目标函数。

现在可以回到开头的数字，理解为什么对数惩罚会那么剧烈。

三个模型对同一个真实标签 \(y = [1, 0, 0]\) 给出预测：

第一个模型比较准确但保留不确定性：\(q = [0.90, 0.05, 0.05]\)，损失 \(-\log(0.90) \approx 0.105\)。

第二个模型犹豫不决：\(q = [0.40, 0.30, 0.30]\)，损失 \(-\log(0.40) \approx 0.916\)。

第三个模型极度错误且自信：\(q = [0.01, 0.98, 0.01]\)，损失 \(-\log(0.01) \approx 4.605\)。

对数让``把概率质量给错类别''比``均匀摇摆''付出更大的代价。正是这个性质驱使模型在不确定时输出诚实的概率——与其把 98\% 的概率押在一个可能错误的类别上，不如均匀分配以降低最坏情况的损失。Log loss 鼓励 calibration 不是来自正则化或先验，而是来自它本身的函数形状。

但这也意味着标签噪声的破坏力被放大。如果训练集中有 1\% 的标签被错误标注，模型可能在那些样本上被惩罚得极重——因为模型被训练去信任标签，而标签是错的。少数高置信错误样本可以显著扭曲决策边界。

\subsection{Cross-entropy 与 KL 之间缺了一项}

常见的说法``cross-entropy 就是 KL''漏掉了 \(H(p)\) 这一项。当 \(p\) 固定不变时——比如同一个训练集上比较两个模型——\(H(p)\) 是常数，省略它不影响最优参数的判定。但跨数据集、跨任务比较损失的绝对值时，这一项不能消失。

假设两个不同的分类任务。任务 A 的标签分布是 \([0.5, 0.3, 0.2]\)，熵 \(H(p_A) \approx 1.49\) nat。任务 B 的标签分布是 \([0.98, 0.01, 0.01]\)，熵 \(H(p_B) \approx 0.11\) nat。两个任务上各自训练到最优，任务 A 的最低可达 loss 约 1.49，任务 B 约 0.11。不能说任务 A 的模型更差——只是任务 A 本身更难，数据固有的不确定性更高。

类似地，one-hot 单标签的 entropy 是零——一个 one-hot 向量没有不确定性。但这不意味着训练损失可以达到零。模型需要从有限样本中推断 \(x\) 到标签的映射关系，类别之间可能有重叠，函数族未必能完全拟合真实条件分布。单样本 target distribution 的熵和总体条件分布的熵是两个层次。

对 \(K\) 类均匀预测 \(q = [1/K, \ldots, 1/K]\)，单样本损失为 \(\log K\)。这是一个有用的 baseline：如果你的训练损失高于 \(\log K\)，模型的表现还不如瞎猜。log 的底决定单位：自然对数给出 nat，以 2 为底给出 bit。换底只把全部损失乘以同一个常数，不改变 minimizer；但数值尺度的改变会影响学习率的选择和早停阈值。

\subsection{数值稳定不是额外工程，是数学目标不被破坏的前提}

Softmax 的输出由 logits \(z\) 通过 \(q_k = e^{z_k} / \sum_j e^{z_j}\) 产生。直接先算 softmax 再取 log 会在极端 logits 上翻车——某个 \(z_k\) 特别大时 \(e^{z_k}\) 溢出，特别小时下溢为零，再取 log 就是 \(-\infty\) 或 NaN。

正确的实现直接把 log-softmax 写成：

\[
\log q(y=k\mid x) = z_k - \operatorname{LSE}(z),
\]

其中 \(\operatorname{LSE}(z) = \log\sum_j e^{z_j}\)，实际计算时先减去 \(\max_j z_j\) 再取指数求和再取 log 再加回最大值。这个等价变换保持了数学目标的每个比特都不变——它只是重新组织了运算顺序以避免浮点灾难。

概率裁剪 \(q_k \leftarrow \max(q_k, \varepsilon)\) 或给 log 输入加一个 \(\varepsilon\) 偏移则不同：它改变了目标函数本身。裁剪后的损失不再是原始 cross-entropy，极小值点也会偏移。如果裁剪不可避免（比如下游计算需要严格正的 \(q_k\)），要清楚它在修改优化目标，而不是在``稳定''原始目标。

\subsection{换了 target 分布，最优解就换了一个}

Label smoothing 把 one-hot target 替换为 \((1-\alpha) \times \text{one-hot} + \alpha \times \text{均匀分布}\)。最小化这个修改版 cross-entropy 的最优 \(q\) 不再等于真实条件分布，而是等于平滑后的 target。如果目标是让模型输出接近真实条件概率，label smoothing 是在有意引入一个 bias——用校准度的代价换取更好的泛化。

Class weighting 通过给不同类别的损失项乘以不同权重，改变了不同标签的相对风险。这在类别不平衡时很实用，但它回答的问题不再是``模型是否接近真实条件分布''，而是``模型在加权风险下是否最优''。

知识蒸馏使用教师模型的软输出作为 target 分布。这时的 cross-entropy 衡量的是学生分布与教师分布的距离，教师本身可能已经偏离了真实数据分布。

所有这些变体仍然可以写成 cross-entropy 的形式，仍然可以用同一套数值稳定技巧，但最优解指向的目标不再相同。如果下游评估关心的是自然人群中的校准度，就必须在未被加权、未被平滑、未被重采样的验证分布上重新检查——训练损失的数字不能直接搬过来当证据。

信息论中的编码基础见 \hyperref[ch:054]{``信息与熵：怎样把``不确定''变成可计算的量''}。
